% ============================================================
\section{Metodologia}\label{sec:metodologia}
% ============================================================

% ------------------------------------------------------------
\subsection{Corpus: da Wikipédia aos parágrafos}\label{sec:dados}

Os vértices são ocorrências de tokens em parágrafos da Wikipédia em português, de oito
\textbf{temas}. Cada tema é o conjunto de artigos alcançados a partir de uma categoria raiz
(Física, Biologia, Economia, Política, Ciência da computação, Música, Futebol e Geografia),
escolhidas distantes entre si e de modo que as palavras ambíguas da P4 tenham sentidos
diferentes em pelo menos dois temas: \emph{banco} é instituição financeira em Economia, banco
de dados em Computação e banco de areia em Geografia. O modelo nunca vê o tema; ele é só um
rótulo da análise.

\paragraph{Busca.} Para cada tema, uma busca em largura desce pelas subcategorias até a
profundidade 2, sem abrir categorias de manutenção nem de listas, e guarda até 600 títulos; os
títulos são embaralhados com uma semente fixa e baixados até completar 150 artigos por tema,
com o identificador da versão baixada de cada artigo registrado. O corte não é aleatório: a API lista as subcategorias em
ordem alfabética, e a busca abriu só parte delas (de 7 de 79 em Computação a 98 de 439 em
Música). Os temas são, portanto, ``artigos próximos da categoria raiz'', e não a área
inteira, o que basta para ter temas bem separados. Como alguns sentidos quase não apareciam no
tema que os representa (\emph{banco} de areia tinha 4 ocorrências em Geografia), 16
\textbf{categorias extras}, cada uma ligada a um tema (Numismática a Economia, Bancos de dados
a Computação, Oceanografia a Geografia\ldots), acrescentam até 60 artigos cada.

\paragraph{Descartes e limpeza.} Um artigo é descartado se for alcançado a partir de mais de
um tema (108 \textbf{páginas multitema}), se for uma biografia (68), se não sobrar nenhum parágrafo de prosa (295), se for uma 
página de desambiguação, que apenas lista os diferentes sentidos de um termo (4) 
ou se estiver duplicado (3). Ficam os parágrafos de prosa corrida (pelo menos 40 palavras, terminados como frase, com
no máximo 8\% de dígitos, sem repetição), divididos em sentenças. Dos \nArtigosCandidatos{}
candidatos, ficaram \nArtigos{} artigos, com \nParagrafos{} parágrafos e \nTokensCorpus{}
tokens do modelo.

% ------------------------------------------------------------
\subsection{Tokens e a amostra de ocorrências}\label{sec:amostra}

O modelo lê \textbf{tokens}, que podem ser palavras inteiras, pedaços de palavra, pontuação
ou dígitos; no Qwen3, o espaço antes de uma palavra faz parte do token (\tok{\esp banco}).
Chamamos de \textbf{tipo} um elemento do vocabulário e de \textbf{ocorrência} uma aparição de
um tipo numa posição de um parágrafo; o número de ocorrências do tipo $t$ na amostra é $f_t$.
Cada ocorrência recebe uma \textbf{categoria de token}: palavra inteira (\tok{\esp banco}),
início de palavra (\tok{\esp torn} em ``tornou''), continuação (\tok{ou} em ``tornou''),
pontuação ou número. 
Quando uma palavra é dividida em vários tokens, apenas o primeiro é classificado como início; todos os demais são continuações. 
Tokens só de espaço não viram vértices. O \emph{tokenizer} do Qwen3 divide muitas palavras do
português em mais de um token (1,76 tokens por palavra, em média). Por isso, os resultados
são apresentados por categoria de token. Na P4, para comparar diferentes ocorrências de uma
palavra, escolhemos um mesmo tipo de token. Por exemplo, para ``banco'', usamos o tipo
\tok{\esp banco}, em minúsculas e com espaço antes; \tok{\esp Banco}, com maiúscula, é outro
tipo. Todas as ocorrências de um mesmo tipo começam com o mesmo vetor, isto é, a mesma lista
de números atribuída ao token pelo modelo. Nas camadas seguintes, esse vetor é transformado
conforme o contexto, permitindo comparar, por exemplo, ``conta no banco'' e ``sentou no banco''.

As ocorrências selecionadas para os estratos de palavras-alvo, controles e multitema
correspondem a palavras inteiras. Na análise da P4, porém, a comparação reúne as ocorrências
pelo tipo de token e pode incluir ocorrências do núcleo em que esse mesmo token é parte de uma
palavra maior, como \tok{\esp capital} em ``capitalismo''.

A amostra limita cada tipo a $f_{\max}=50$ ocorrências, usa uma semente fixa e tem quatro
\textbf{estratos}:
\begin{enumerate}
    \item \textbf{Núcleo} (P1 a P3): 5 artigos por tema e 2 parágrafos de cada, 80 parágrafos
    inteiros. Os tokens deles viram vértices, inclusive pedaços de palavra, números e
    pontuação (só os espaços soltos ficam de fora), até o limite $f_{\max}$: os 25 tipos mais
    frequentes, como a vírgula, \tok{\esp de} e o ponto, ficam com 50 ocorrências sorteadas, o
    que corta 1.915 dos 11.935 tokens e deixa \nVerticesNucleo{} vértices. É o único estrato
    com parágrafos inteiros.
    \item \textbf{Palavras-alvo} (P4): palavras com sentidos diferentes em temas diferentes,
    com uma cota de 16 a 25 ocorrências em cada tema de sentido, no máximo uma por parágrafo e
    três por artigo. Ficaram \nAlvosMantidos{} (\emph{banco}, \emph{campo}, \emph{órgão},
    \emph{estado}, \emph{carga}, \emph{rede}, \emph{nota}, \emph{capital}, \emph{família},
    \emph{massa} e \emph{matriz}), com \nVerticesAlvo{} vértices; \emph{planta} saiu por falta
    de ocorrências.
    \item \textbf{Controles} (P4): seis palavras de sentido estável (\emph{ano},
    \emph{século}, \emph{grande}, \emph{primeiro}, \emph{importante}, \emph{água}), com a mesma
    alocação de cotas dos alvos, para que uma diferença entre os grupos não venha do desenho:
    \nVerticesControle{} vértices.
    \item \textbf{Multitema} (P4): 150 palavras de conteúdo que aparecem em pelo menos três
    temas, como \emph{metal}, \emph{luz}, \emph{guerra} e \emph{motor}, com até 20 ocorrências
    de cada. Elas permitem aplicar a análise também a palavras selecionadas por aparecerem em
    vários temas, independentemente de terem sido escolhidas por sua ambiguidade:
    \nVerticesMultitema{} vértices.
\end{enumerate}
Nos estratos 2 a 4, o parágrafo da ocorrência passa pelo modelo só como contexto: só as
ocorrências sorteadas viram vértices. Em geral há uma por parágrafo; quando o sorteio escolhe
palavras diferentes do mesmo parágrafo, o parágrafo entra uma vez só, com todas elas (388 dos 3.150
parágrafos desses estratos têm de 2 a 5 vértices). Ao todo são \nVertices{} vértices de \nTipos{} tipos, em
\nSequencias{} parágrafos de 1.172 artigos.

\paragraph{Anotação dos sentidos.}\label{sec:dados-sentidos} O tema é só um indicador do
sentido: um artigo de Computação pode citar o Banco Central. Por isso, as 508 ocorrências das
palavras-alvo foram classificadas pelo modelo de linguagem Claude Opus~5.5, sem revisão humana.
Para cada palavra, foram definidos de três a sete rótulos de sentido; cada ocorrência recebeu
um deles. Por exemplo, as ocorrências de \emph{órgão} foram classificadas como biológicas,
institucionais ou referentes ao instrumento musical. Cada ocorrência foi lida com o trecho do
parágrafo, sem o tema nem o título do artigo. Como o modelo
analisado só vê o texto à esquerda, cada ocorrência recebeu também uma marca que diz se esse
trecho já decide o sentido (442 das 508).

% ------------------------------------------------------------
\subsection{Modelo e vetores}\label{sec:modelo}

O modelo é o \textbf{Qwen3-4B-Base}~\cite{qwen3}, um Transformer de 4 bilhões de parâmetros,
com 36 blocos e vetores de dimensão 2560, numa revisão fixa e em \emph{bfloat16}. A versão
\emph{Base} é treinada só para prever o próximo token. A matriz de \emph{embeddings} $E$ da
entrada é a mesma da saída~\cite{press2017}: a linha $e_t$ de um token é o vetor com que ele
entra no modelo e também a direção usada para prevê-lo.

Cada parágrafo passa \textbf{sozinho} pelo modelo, precedido do separador
\tok{<|endoftext|>}, sem o título nem os outros parágrafos do artigo. Como a atenção é causal,
o vetor de uma posição depende só do seu token e dos anteriores. Cada ocorrência $i$, de tipo
$t_i$, recebe quatro vetores: \cod{lex}, a linha $e_{t_i}$, igual para todas as ocorrências do
tipo; e \cod{L01}, \cod{L18} e \cod{L36}, a saída bruta dos blocos 1, 18 e 36 na posição de
$i$ (antes da normalização final), lida por \emph{hooks} do PyTorch. Uma segunda execução, com
a mesma amostra, guarda a saída dos 36 blocos, para as curvas camada a camada. Guardamos também
a saída normalizada do bloco 36 (\cod{L36n}) e o token que o modelo prevê depois de cada
ocorrência. 

Quando duas ou mais ocorrências têm exatamente a mesma sequência de tokens desde o
início do parágrafo até a própria ocorrência, o modelo causal deveria produzir vetores iguais
para elas. Para evitar que pequenas diferenças de arredondamento criem distinções artificiais,
copiamos o vetor da primeira ocorrência para as demais.

% ------------------------------------------------------------
\subsection{Construção das redes e o tratamento dos empates}\label{sec:knn}

A semelhança entre dois vértices é o cosseno $s_{ij}=x_i^\top x_j/(\lVert x_i\rVert\,\lVert
x_j\rVert)$. Cada vértice $i$ aponta para os $k=10$ vértices de
maior cosseno com ele, entre todos os outros vértices da amostra, de qualquer parágrafo, artigo
ou tema; esse conjunto é a vizinhança $\N_i$. A rede é \textbf{direcionada} ($i\to j$ não implica
$j\to i$) e sem pesos; todo vértice tem grau de saída 10. As quatro redes diferem só nos arcos
(Figura~\ref{fig:ilustracao}). Usamos a versão \textbf{direcionada} nas medidas de vizinhança,
no grau de entrada e na reciprocidade; a \textbf{não direcionada por união} (aresta se $i$
escolhe $j$ ou $j$ escolhe $i$) na clusterização, nas distâncias, nos componentes e nas
comunidades.

\begin{figure}[htbp]
\centering
\begin{tikzpicture}[
    scale=0.86, every node/.style={font=\small},
    v/.style={draw, thick, rounded corners=3pt, inner sep=2.5pt, minimum height=5.5mm,
        font=\footnotesize},
    eco/.style={v, fill=green!22},
    com/.style={v, fill=magenta!18},
    geo/.style={v, fill=orange!28},
    arco/.style={-{Stealth[length=2.2mm]}, thick},
    fora/.style={-{Stealth[length=2mm]}, thick, dashed, black!45},
    titulo/.style={font=\small\bfseries},
]
% (a) rede lexical
\begin{scope}
    \node[titulo] at (0,3.0) {(a) rede lexical (\texttt{lex})};
    \node[eco] (bE) at (0,1.0) {banco};
    \node[com] (bC) at (-1.3,-0.5) {banco};
    \node[geo] (bG) at (1.3,-0.5) {banco};
    \node[eco] (j) at (0,2.2) {juros};
    \node[com] (d) at (-2.7,-1.5) {dados};
    \node[geo] (a) at (2.7,-1.5) {areia};
    \draw[arco] (bE) to[bend right=10] (bC);
    \draw[arco] (bC) to[bend right=10] (bE);
    \draw[arco] (bE) to[bend left=10] (bG);
    \draw[arco] (bG) to[bend left=10] (bE);
    \draw[arco] (bC) to[bend right=10] (bG);
    \draw[arco] (bG) to[bend right=10] (bC);
    \draw[fora] (j.east) -- ++(0.9,0.2);
    \draw[fora] (j.west) -- ++(-0.9,0.2);
    \draw[fora] (d.south) -- ++(-0.4,-0.6);
    \draw[fora] (d.south) -- ++(0.4,-0.6);
    \draw[fora] (a.south) -- ++(0.4,-0.6);
    \draw[fora] (a.south) -- ++(-0.4,-0.6);
\end{scope}
% (b) rede contextual
\begin{scope}[xshift=8.2cm]
    \node[titulo] at (0,3.0) {(b) rede de uma camada contextual};
    \node[eco] (bE) at (0,1.0) {banco};
    \node[com] (bC) at (-1.3,-0.5) {banco};
    \node[geo] (bG) at (1.3,-0.5) {banco};
    \node[eco] (j) at (0,2.2) {juros};
    \node[com] (d) at (-2.7,-1.5) {dados};
    \node[geo] (a) at (2.7,-1.5) {areia};
    \draw[arco] (bE) to[bend left=12] (j);
    \draw[arco] (j) to[bend left=12] (bE);
    \draw[arco] (bC) to[bend left=12] (d);
    \draw[arco] (d) to[bend left=12] (bC);
    \draw[arco] (bG) to[bend right=12] (a);
    \draw[arco] (a) to[bend right=12] (bG);
    \draw[arco] (bE) -- (bC);
    \draw[arco] (bC) -- (bG);
    \draw[arco] (bG) -- (bE);
    \draw[fora] (j.east) -- ++(0.9,0.2);
    \draw[fora] (d.south) -- ++(0,-0.65);
    \draw[fora] (a.south) -- ++(0,-0.65);
\end{scope}
\end{tikzpicture}
\caption{Exemplo ilustrativo (não são dados reais), com $k=2$: seis ocorrências, \emph{banco}
em três parágrafos de temas diferentes (verde: Economia; rosa: Computação; laranja: Geografia)
e uma palavra de cada um desses parágrafos; as setas tracejadas vão para vértices fora do
desenho. (a) Na rede lexical, as três ocorrências de \emph{banco} têm o mesmo vetor e apontam
umas para as outras. (b) Numa camada contextual, cada \emph{banco} passa a apontar também para
uma palavra do seu contexto. Os vértices são os mesmos; só os arcos mudam.}\label{fig:ilustracao}
\end{figure}

\paragraph{Empates.} Para escolher os vizinhos de cada palavra, o método procura as ocorrências
mais parecidas com ela. Quando há mais candidatos igualmente parecidos do que vagas disponíveis,
as vagas restantes são preenchidas por sorteio.

Isso acontece com frequência na rede lexical, pois diferentes ocorrências de uma mesma palavra
têm a mesma representação. O sorteio evita que as mesmas ocorrências sejam sempre escolhidas
apenas por aparecerem primeiro na lista. O sorteio é repetido com 20 sementes em todas as
representações. Nas medidas de vizinhança, usam-se todas elas. Nas métricas estruturais, a
média das 20 sementes é usada apenas para a rede lexical; nas camadas contextuais, usa-se uma
semente, pois os sorteios quase não mudam a rede.

\paragraph{A rede do vocabulário.} Além das redes construídas com as palavras da amostra,
criamos uma rede que reúne todo o vocabulário do modelo. Nela, cada vértice representa um token
diferente, e cada token é ligado aos 10 tokens com representações mais parecidas com a sua.

A rede tem 151.643 vértices: inclui todos os tokens do vocabulário, exceto os 26 tokens especiais.

% ------------------------------------------------------------
\subsection{Medidas}\label{sec:medidas}

\paragraph{Métricas estruturais.} Calculadas com o igraph~\cite{csardi2006}: vértices,
arestas, grau médio, densidade, reciprocidade, clusterização global (transitividade) e local
média, componentes, distância média e diâmetro (exatos, no maior componente), e a distribuição
do grau de entrada, cuja assimetria (\emph{skewness}) mede a \emph{hubness}~\cite{radovanovic2010}.

\paragraph{Troca de vizinhos e piso de ruído (P1, P2).} A vizinhança de uma ocorrência nas
redes $a$ e $b$ é comparada pelo índice de Jaccard $J_i(a,b)=|\N_i^{(a)}\cap\N_i^{(b)}|/
|\N_i^{(a)}\cup\N_i^{(b)}|$; com 10 vizinhos e $c$ em comum, $J_i=c/(20-c)$ ($J=0{,}11$
quando sobram 2 vizinhos). Na rede lexical, duas sementes do sorteio já dão $J<1$ sem mudança
nenhuma de representação. O \textbf{piso de ruído} $J^{\text{ruído}}_i$ é o Jaccard médio entre
redes lexicais de sementes diferentes; ele depende da frequência (0,12 para tipos com 50
ocorrências, 0,39 de 10 a 49, cerca de 0,75 para tipos raros). A \textbf{mudança acima do
piso} é $M^{+}_i=J^{\text{ruído}}_i-J_i(\mathrm{lex},\ell)$, a parte da troca que o sorteio não
explica. O Jaccard por composição $J^w$ compara quantos vizinhos de cada \emph{tipo} a
ocorrência tem nas duas redes e quase não depende do sorteio.

\paragraph{Dominância lexical (P1, P2).} É a fração dos vizinhos que são o mesmo tipo,
$D_i=|\{j\in\N_i:t_j=t_i\}|/k$. Na rede lexical ela atinge o teto $\min(f_{t_i}-1,k)/k$, e a
dominância relativa ao teto $\tilde D_i$ permite comparar tipos de frequências diferentes. O
nível do acaso é $(f_{t_i}-1)/(n-1)\le0{,}004$.

\paragraph{Comunidades e concordância (P3).} As comunidades são encontradas na rede por união
com o Leiden~\cite{traag2019}, que maximiza a modularidade~\cite{newman2006} (10 execuções; a
concordância média entre elas é a \emph{estabilidade}). A partição é comparada com rótulos dos
vértices (tipo de token, tema, artigo, parágrafo, sentença, próximo token, token previsto pelo
modelo e faixa de posição) pela informação mútua normalizada
(NMI), só nos vértices do núcleo: as comunidades vêm da rede inteira, mas só o núcleo tem
parágrafos inteiros; nos outros estratos, quase nove em cada dez parágrafos têm um só vértice. Como o NMI cresce com o número de classes mesmo sem relação nenhuma~\cite{vinh2010}, e
os rótulos vão de 4 a 2.780 classes, cada NMI é descontado de uma \textbf{linha de base}, o NMI
médio com o rótulo embaralhado 20 vezes; a medida lida é \textbf{NMI $-$ base}.

\paragraph{Separação das ocorrências (P4).} Para cada tipo com pelo menos 10 ocorrências, em
quatro grupos (alvos, controles, multitema e palavras funcionais), medimos a
\textbf{autossimilaridade} (cosseno médio entre as ocorrências~\cite{ethayarajh2019}), a
diferença \textbf{intra $-$ entre} (cosseno médio entre ocorrências do mesmo tema menos o
cosseno entre temas diferentes; positiva se elas se agrupam por tema), o \textbf{número efetivo
de comunidades} ocupadas ($e^{H}$, com $H$ a entropia da distribuição pelas comunidades) e o
\textbf{NMI entre comunidade e tema} dentro do tipo, com a sua linha de base. Nos alvos, as
medidas são repetidas com o sentido anotado no lugar do tema.

\paragraph{Vizinhos no vocabulário.} Como $E$ é usada na entrada e na saída, ordenamos as
151.669 linhas de $E$ pelo cosseno com o vetor de cada ocorrência, em cada bloco, e anotamos a
posição do \emph{próprio} token e a do \emph{próximo} token do texto, uma versão simples da
``lente de \emph{logits}''~\cite{nostalgebraist2020}. Na saída normalizada do bloco 36, a linha
de maior produto interno coincide com a previsão do modelo em \concordanciaPredicao{} das
ocorrências, o que valida o cálculo.
